\scanpage{85}
\noindent Définition \uline{directe} de l'extension $\underline{\mathcal M}_{X/S}$.

Soit $\mathcal M^\sharp_{X/S}$ le foncteur $(\mathrm{Sch})_{/S}^{\circ}\longrightarrow(\mathrm{Ens})$, faisceau \inserted{fppf} associé au préfaisceau
\[
\begin{aligned}
 &\deleted{\mathcal M'_{X/S}(T)=}\\
 Q:\quad T&\rightsquigarrow \text{classes, à isom près, de faisceaux}\\
 &\phantom{\rightsquigarrow{}}\text{inversibles sur }X\times_S T,\text{ munis d'une}\\
 &\phantom{\rightsquigarrow{}}\text{\uline{connexion} relativement à }T.
\end{aligned}
\]

\deleted{On a une} \deleted{N.B} On a une suite exacte de pré\uline{faisceaux} en $T$
\[
 H^0(X,\underline O_X^*)\xrightarrow{\varphi\mapsto\frac{d\varphi}{\varphi}}
 \frac{H^0(X,\Omega^1_X)}{\deleted{\unsure{\operatorname{Im}|d\log}}}
 \longrightarrow Q(T)\longrightarrow\mathrm{Pic}(X_T)\longrightarrow H^1(X,\underline\Omega^1_{X/S})
\]
d'où en passant aux faisceaux en $T$ associés
\[
 0\longrightarrow f_*(\underline\Omega^1_X)/\underbrace{\operatorname{Im}f_*(\underline O_X^*)}_{\text{par }\varphi\mapsto\frac{d\varphi}{\varphi}}
 \longrightarrow\mathcal M^\sharp_{X/S}\longrightarrow\underline{\mathrm{Pic}}^{w}_{X/S}\longrightarrow0
\]
où on a défini encore
\[
 \underline{\mathrm{Pic}}^w_{X/S}=\operatorname{Ker}\bigl(\underline{\mathrm{Pic}}_{X/S}\longrightarrow R^1 f_*(\underline\Omega^1_{X/S})\bigr)
\]

On va définir un homomorphisme d'ext.
\diagram{\begin{tikzcd}[ampersand replacement=\&,column sep=small]
0\arrow[r,dash]\&\deleted{\unsure{f_*|f}}\underline{\unsure{\omega|\Omega}}_{X/S}\arrow[r]\arrow[d,"{\mathrm{can}}"]\& f_*(\underline E_{X/S})\arrow[r,dash]\arrow[d,"\alpha"]\&\underline{\mathrm{Div}}^w_{X/S}\arrow[r,dash]\arrow[d,"{\mathrm{can}}"]\&0\\
0\arrow[r]\&\underline\omega_{X/S}/\operatorname{Im}f_*(\underline O_X^*)\arrow[r]\&\mathcal M^\sharp_{X/S}\arrow[r]\&\underline{\mathrm{Pic}}^w_{X/S}\arrow[r]\&0
\end{tikzcd}}

Pour ceci, on va interpréter \deleted{$\underline E_{X/S}$} l'extension $\underline E_{X/S}$ de $\underline{\mathrm{Div}}_{X/S}$ par $\underline\Omega^1_{X/S}$, en notant que pour un diviseur relatif donné $D$ sur un
% archive sheet number: 84

